% !TeX root = ../se200_identity_regimes.tex

\subsection{Refinement of Enduring Non-Normative Referents}
\label{subsec:enr_split}

This subsection shows that the class of enduring non-normative referents
does not define a single regime profile under the transformation basis,
and must be refined.

%--------------------------------------------
\subsubsection{Initial ENR Profile}

\begin{definition}[Provisional ENR profile]
A provisional ENR profile is a regime profile whose identity carrier is an
enduring non-normative referent, i.e., an enduring entity that does not itself
ground obligations or normative structure.
\end{definition}


Under this provisional formulation, identity is intended to track persistence
of an enduring referent independently of normative or causal interpretation.


%--------------------------------------------
\subsubsection{Transformation Behavior}

We examine the classification of transformations under the provisional ENR profile.

\begin{itemize}
\item $\mathrm{RE}$: IGN
\item $\mathrm{AN}$: IGN
\item $\mathrm{RF}$: PRS
\item $\mathrm{AD}$: PRS
\item $\mathrm{RC}$: IGN
\item $\mathrm{RA}$: N/A
\item $\mathrm{SU}$: BRK
\item $\mathrm{BF}$: PRS or BRK
\item $\mathrm{PV}$: IGN
\item $\mathrm{SE}$: PRS
\end{itemize}


The classification of branching is indeterminate under the provisional ENR profile:
it is identity-preserving for locus-bound referents but identity-breaking
for instrument-bound referents.
This indeterminacy is the source of split pressure.


%--------------------------------------------
\subsubsection{Split Pressure}

\begin{proposition}[ENR split pressure]
\label{prop:enr_split_pressure}
The provisional ENR profile is under split pressure with respect to
the branching transformation $\mathrm{BF}$.
\end{proposition}

\begin{proof}[Proof of Proposition~\ref{prop:enr_split_pressure}]
Let $x \in \mathcal{R}$ be a locus-bound representation and let $\tau$ be a
branching transformation producing continuations over the same locus.
Since the locus persists through branching,
$\mathrm{Class}_P(\mathrm{BF}) = \mathrm{PRS}$ is required,
i.e., branching must be identity-preserving.

Let $x' \in \mathcal{R}$ be an instrument-bound representation and let $\tau'$ be a
branching transformation producing two non-identical artifact continuations.
At most one continuation can be the same artifact as $x'$,
so $\mathrm{Class}_P(\mathrm{BF}) = \mathrm{BRK}$ is required,
i.e., branching must be identity-breaking.

These requirements are contradictory.
No single value of $\mathrm{Class}_P(\mathrm{BF})$ is consistent with both classes of cases.
Therefore the provisional ENR profile is under split pressure with respect to $\mathrm{BF}$.
\end{proof}

[Informally]
After branching, the locus remains the same but the artifact does not.


%--------------------------------------------
\subsubsection{Separated Profiles}

We therefore define two distinct regime profiles.

\begin{definition}[ENR-L profile]
The locus-bound profile $\mathrm{ENR\mbox{-}L}$ is defined by:
\begin{itemize}
\item carrier: enduring locus,
\item identity basis: persistence of the same locus,
\item transformation classification:
\begin{itemize}
\item $\mathrm{RE}$: IGN
\item $\mathrm{AN}$: IGN
\item $\mathrm{RF}$: PRS
\item $\mathrm{AD}$: PRS
\item $\mathrm{RC}$: IGN
\item $\mathrm{RA}$: N/A
\item $\mathrm{SU}$: BRK
\item $\mathrm{BF}$: PRS
\item $\mathrm{PV}$: IGN
\item $\mathrm{SE}$: PRS
\end{itemize}
\end{itemize}
\end{definition}

\begin{definition}[ENR-I profile]
The instrument-bound profile $\mathrm{ENR\mbox{-}I}$ is defined by:
\begin{itemize}
\item carrier: enduring artifact or instrument,
\item identity basis: persistence of the same artifact,
\item transformation classification:
\begin{itemize}
\item $\mathrm{RE}$: IGN
\item $\mathrm{AN}$: IGN
\item $\mathrm{RF}$: PRS
\item $\mathrm{AD}$: PRS
\item $\mathrm{RC}$: IGN
\item $\mathrm{RA}$: N/A
\item $\mathrm{SU}$: BRK
\item $\mathrm{BF}$: BRK
\item $\mathrm{PV}$: IGN
\item $\mathrm{SE}$: PRS
\end{itemize}
\end{itemize}
\end{definition}

%--------------------------------------------
\subsubsection{Non-Collapse}

\begin{proposition}[ENR-L and ENR-I do not collapse]
\label{prop:enr_noncollapse}
The profiles $\mathrm{ENR\mbox{-}L}$ and $\mathrm{ENR\mbox{-}I}$ induce
distinct equivalence relations on $\mathcal{R}$.
\end{proposition}

\begin{proof}[Proof of Proposition~\ref{prop:enr_noncollapse}]
Let $x \in \mathcal{R}$ be an instrument representation.
Let $\tau$ be a branching transformation producing a distinct artifact
continuation $y := \tau(x)$ over the same locus as $x$.

Under $\mathrm{ENR\mbox{-}L}$:
\[
\mathrm{Class}_{\mathrm{ENR\mbox{-}L}}(\mathrm{BF}) = \mathrm{PRS},
\]
so $\tau$ is identity-preserving and $x \sim_{\mathrm{ENR\mbox{-}L}} y$.

Under $\mathrm{ENR\mbox{-}I}$:
\[
\mathrm{Class}_{\mathrm{ENR\mbox{-}I}}(\mathrm{BF}) = \mathrm{BRK},
\]
so $\tau$ is identity-breaking and $x \not\sim_{\mathrm{ENR\mbox{-}I}} y$.

Therefore $\sim_{\mathrm{ENR\mbox{-}L}} \neq \sim_{\mathrm{ENR\mbox{-}I}}$.
\end{proof}

[Informally]
After branching, two representations can remain identical as loci
even when they are no longer the same artifact.
Thus the profiles define different equivalence classes.


%--------------------------------------------
\subsubsection{Result}

\begin{theorem}[ENR refinement]
The class of enduring non-normative referents must be refined into exactly two
distinct regime profiles under the branching transformation $\mathrm{BF}$:
\[
\mathrm{ENR\mbox{-}L}
\quad\text{and}\quad
\mathrm{ENR\mbox{-}I}.
\]
\end{theorem}

[Informally]
Two representations can refer to the same locus even when the artifacts differ,
but cannot refer to the same artifact once branching produces distinct objects.

Exactly two profiles suffice: $\mathrm{BF}$ admits only $\mathrm{PRS}$
or $\mathrm{BRK}$, and each is realized by one of the two profiles.
